% Ensure included PDFs (banners) with newer versions embed without warnings
\pdfminorversion=7
\documentclass[aspectratio=169]{beamer}

\makeatletter
% Make LaTeX find theme .sty files in Template
\def\input@path{{Template/}}
\makeatother
\usepackage{graphicx}
% Make graphics (including banner PDFs) resolvable without TEXINPUTS
\graphicspath{{Template/}{Template/banners/}{../Common_Images/}}

%================================================================%
% Theme and Package Setup
%================================================================%
\usetheme{ucl}
\setbeamercolor{banner}{bg=darkpurple}

% Navigation and Footer
\setbeamertemplate{navigation symbols}{\vspace{-2ex}}
\setbeamertemplate{footline}[author title date]
\setbeamertemplate{slide counter}[framenumber/totalframenumber]

\usepackage[utf8]{inputenc}
\usepackage[british]{babel}
\usepackage{amsmath, amssymb, amsthm}
\usepackage{tikz}
\usepackage{pgfplots}
\pgfplotsset{compat=1.17}
\usetikzlibrary{calc, positioning, arrows.meta, shapes.geometric, trees, backgrounds, shapes.misc, graphs, quotes, shadows, fit, decorations.pathreplacing} 
\usefonttheme{professionalfonts}
\usepackage{eulervm}
\usepackage{listings}
\usepackage{xcolor}
\usepackage{algorithm}
\usepackage{algpseudocode}
\usepackage[square,authoryear]{natbib}

% Define custom colours
\makeatletter
\@ifundefined{color@stone}{%
    \definecolor{stone}{gray}{0.95}%
}{}
\makeatother
\definecolor{darkgreen}{rgb}{0.0, 0.5, 0.0}
\definecolor{uclblue}{cmyk}{1,0.24,0,0.64}

\setbeamercovered{transparent}

% Define theoremblock
\makeatletter
\newenvironment<>{theoremblock}[1]{%
    \begin{block}#2{#1}%
}{\end{block}}
\makeatother

% Section divider slides
\AtBeginSection[]{
  \begin{frame}
    \frametitle{\textbf{\Large\insertsectionhead}}
    \begin{center}
        \huge\textbf{\insertsectionhead}
    \end{center}
  \end{frame}
}

% Operator definitions
\DeclareMathOperator*{\argmin}{arg\,min}
\DeclareMathOperator{\prox}{prox}
\DeclareMathOperator{\proj}{proj}
\DeclareMathOperator{\divg}{div}

\bibliographystyle{plainnat}

\title{Deblurring Challenge}
\subtitle{Restoring Images from Randomised Motion Blur}
\author[Martin Benning (University College London)]{Martin Benning}
\date[CCMI]{CCMI -- Deep Learning in Theory and Practice \\[1cm] Monday, 18 May 2026}
\institute[]{University College London}

\begin{document}

% --- Title Frame ---
\begin{frame}
  \titlepage
\end{frame}

%================================================================%
\section{The Deblurring Challenge}
%================================================================%

\begin{frame}{Overview of the Challenge}
    \textbf{The Task:} Develop a data-driven model to restore high-quality images from degraded measurements suffering from \textbf{randomised motion blur} and \textbf{noise}.
    
    \vspace{1em}
    \begin{columns}[T]
        \column{0.5\textwidth}
        \textbf{Key Details:}
        \begin{itemize}
            \item \textbf{Tools:} PyTorch and the DeepInverse (`deepinv`) library.
            \item \textbf{Evaluation:} Hidden test set (PSNR and SSIM).
            \item \textbf{Format:} Friendly competition among PhD student groups via an automated GitHub Actions leaderboard.
        \end{itemize}
        
        \column{0.5\textwidth}
        \begin{block}{The Forward Model}
            We observe
            $$ y^\delta = A x^\dagger + \eta $$
            where $A$ represents a convolution with motion blur kernel, and $\eta$ is measurement noise.
        \end{block}
    \end{columns}
\end{frame}

\begin{frame}{Data Generation \& Training Strategy}
    To train your data-driven models, you will need sharp ground-truth images to generate synthetic degraded data.

    \begin{itemize}
        \item \textbf{Datasets:} Use `torchvision` (e.g., Food101, Flowers102) or DeepInverse datasets (e.g., BSDS500, DIV2K) as your source of sharp images.
        \item \textbf{Synthetic Degradation:} Apply randomised motion blur operators and noise using `deepinv`.
        \begin{itemize}
            \item \textit{Tip:} Blur a large image segment \emph{before} extracting patches to handle boundary conditions correctly.
        \end{itemize}
        \item \textbf{Training Strategy:} Train on small patches (e.g., $128 \times 128$ or $256 \times 256$) to improve efficiency and allow the network to generalise to arbitrary resolutions during inference.
    \end{itemize}
\end{frame}

%================================================================%
\section{The Mathematical Baseline}
%================================================================%

\begin{frame}{Why do we need a Baseline?}
    Before unleashing a deep neural network on an inverse problem, it is critical to understand the classical mathematical approach. This provides a performance lower-bound and ensures our data-driven model is actually learning something useful.
    
    \vspace{0.5em}
    For the baseline, we solve the inverse problem $Ax = y^\delta$ using a \textbf{Variational Method}. We formulate an optimisation problem:
    \begin{align*}
        \hat{x} = \argmin_x \left\{ \frac{1}{2} \| Ax - y^\delta \|^2 + \alpha R(x) \right\}
    \end{align*}
    \begin{itemize}
        \item \textbf{Data Fidelity:} $\frac{1}{2} \| Ax - y^\delta \|^2$ ensures the solution matches the measurements.
        \item \textbf{Regularisation:} $R(x)$ enforces prior knowledge, stabilising the inversion. $\alpha > 0$ balances the two.
    \end{itemize}
\end{frame}

\begin{frame}{Total Variation (TV) Regularisation}
    A powerful and classical choice for imaging is \textbf{Total Variation (TV)}, giving rise to the ROF (Rudin, Osher, Fatemi) model.
    
    \begin{block}{Discrete Total Variation}
        For a 2D image $x$, the isotropic Total Variation is the $\ell_1$-norm of the gradient magnitudes:
        \begin{align*}
            TV(x) = \sum_{i,j} \sqrt{ (\nabla_h x)_{i,j}^2 + (\nabla_v x)_{i,j}^2 }
        \end{align*}
        where $\nabla_h$ and $\nabla_v$ are horizontal and vertical finite differences.
    \end{block}
    
    \textbf{Why TV?} It heavily penalises noise and oscillations, but \emph{preserves sharp edges} (unlike $\ell_2$ smoothness priors which blur edges).
\end{frame}

\begin{frame}{Solving the Baseline: The Condat-V\~{u} Algorithm}
    The objective function is
    \begin{align*}
        \min_x \underbrace{\frac{1}{2} \| Ax - y^\delta \|^2}_{\text{Smooth (Gradient desc.)}} + \underbrace{\alpha TV(x)}_{\text{Non-smooth (Proximal operator)}}
    \end{align*}
    Because $TV(x)$ involves a linear operator ($\nabla$) inside a non-smooth norm, standard gradient descent fails.
    
    \vspace{0.5em}
    Instead, we use \textbf{Primal-Dual Splitting} (PDHG/Chambolle--Pock), specifically the \textbf{Condat-V\~{u} extension} \citep{Pock2009PDHG, Chambolle2011PDHG, condat2013primal, vu2013splitting}. It splits the problem into:
    \begin{enumerate}
        \item A primal gradient step on the data fidelity.
        \item A dual proximal step on the regularisation.
    \end{enumerate}
\end{frame}

\begin{frame}{Condat-V\~{u} Iterations for Deblurring}
    \only<1-2>{Let $p$ be a dual variable (a vector field corresponding to the image gradients). We alternate updates with step sizes $\tau$ (primal) and $\sigma$ (dual):}
    
    \begin{theoremblock}{Algorithm Steps}
        For $k = 0, 1, 2, \dots$ until convergence:
        \begin{align*}
            1.\; \tilde{x}^{k+1} &= x^k - \tau A^\top(A x^k - y^\delta) + \tau \operatorname{div}(p^k) \\
            2.\; x^{k+1} &= \operatorname{proj}_{[0, 1]}(\tilde{x}^{k+1}) \quad \text{\footnotesize(Optional: constrain to valid pixel range)} \\
            3.\; p^{k+1} &= \operatorname{proj}_{B_\alpha}\left( p^k + \sigma \nabla (2 x^{k+1} - x^k) \right)
        \end{align*}
    \end{theoremblock}
    
    \begin{itemize}
        \item<1-> $\operatorname{div} = -\nabla^\top$ is the discrete divergence operator.
        \item<2-> $\operatorname{proj}_{B_\alpha}$ projects the dual variable pointwise onto an $\ell_2$-ball of radius $\alpha$.
        \item<3-> \textbf{Convergence condition:} $\tau \left( \frac{\|A\|^2}{2} + \sigma \|\nabla\|^2 \right) \leq 1$.
    \end{itemize}
\end{frame}

%================================================================%
\section{Rules, Constraints \& Logistics}
%================================================================%

\begin{frame}{Registration \& Leaderboard}
    \textbf{GitHub Classroom Sign-up:}
    \begin{itemize}
        \item Join the assignment via this link: \url{https://classroom.github.com/a/dy4H1JHT}
        \item Sign up in groups of \textbf{ideally 3} (or 4 if necessary). Note that a maximum of 7 teams is allowed in total.
        \item Come up with a \textbf{unique team name} that clearly identifies your group.
        \item The first team member to sign up creates the team and sets the name; the remaining members simply join that established team.
    \end{itemize}
    
    \textbf{Tracking Progress:}
    \begin{itemize}
        \item There will be a live leaderboard to keep track of everyone's progress on the hidden test set.
        \item The leaderboard can be found at: \url{https://deep-learning-theory-practice-2026.github.io/Coding-Challenge-Leaderboard/}
    \end{itemize}
\end{frame}

\begin{frame}[fragile]{Repository Structure}
    \begin{columns}
        \column{0.45\textwidth}
        \begin{lstlisting}[basicstyle=\tiny\ttfamily, language=bash, breaklines=true]
.
|-- .github
|   |-- scripts
|   |   |-- evaluate.py
|   |   `-- unpack_pt.py
|   `-- workflows
|       `-- evaluation.yml
|-- data
|   `-- .gitkeep
|-- .gitignore
|-- environment.yml
|-- LICENSE
|-- model.py
|-- predict.py
|-- requirements.txt
`-- README.md
        \end{lstlisting}
        
        \column{0.55\textwidth}
        \textbf{What you need to edit/provide:}
        \begin{itemize}
            \item \texttt{model.py}: Define your PyTorch network architecture.
            \item \texttt{predict.py}: Implement the inference logic.
            \item \texttt{requirements.txt}: Add any extra Python libraries.
            \item \texttt{*.pth} file: Upload your trained model weights via a GitHub Release.
        \end{itemize}
    \end{columns}
\end{frame}

\begin{frame}{Constraints \& Requirements}
    \begin{alertblock}{Strict Rules for the Auto-Grader}
        \begin{itemize}
            \item \textbf{Model Size:} Your submitted \texttt{.pth} weight file must be $\le$ \textbf{100 MB}.
            \item \textbf{Compute Limits:} Inference is evaluated on CPU only. Your pipeline must process the entire hidden test set in under \textbf{20 minutes}.
            \item \textbf{Dynamic Sizing:} Avoid hardcoding dimensions! Your network must handle variable input image sizes.
        \end{itemize}
    \end{alertblock}
    
    \vspace{0.5em}
    \textbf{Important Warnings:}
    \begin{itemize}
        \item Do \textbf{NOT} modify any files in the \texttt{.github/} directory.
        \item Do \textbf{NOT} alter the command-line arguments (argparser) in \texttt{predict.py}. The auto-grader depends on them.
    \end{itemize}
\end{frame}

\begin{frame}{Testing and Submission Workflow}
    \begin{center}
    \begin{tikzpicture}[
        box/.style={draw, thick, fill=uclblue!20, minimum width=3cm, minimum height=1cm, align=center, rounded corners},
        arr/.style={->, thick, >=latex}
    ]
        \node[box] (train) at (0,0) {1. Train Locally};
        \node[box] (test) at (5,0) {2. Test \texttt{predict.py}};
        \node[box] (push) at (10,0) {3. Push Code \&\\Create Release};
        \node[box, fill=darkgreen!20, minimum height=1.5cm] (grade) at (10,-2.5) {4. GitHub Action\\Evaluates \& Scores};
        
        \draw[arr] (train) -- (test);
        \draw[arr] (test) -- (push);
        \draw[arr] (push) -- (grade);
    \end{tikzpicture}
    \end{center}
    \vspace{0.5em}
    \textbf{Local Testing:}
    We provided \texttt{data/test\_dataset.pt}. Test your setup via:
    \texttt{python predict.py --input data/test\_dataset.pt --output sub\_out}

    \vspace{0.5em}
    \textbf{Submission:} Publishing a new GitHub Release triggers the evaluation. The latest code on \texttt{main} is also evaluated on a schedule before each leaderboard update.
\end{frame}

\begin{frame}{Friday Presentations}
    \textbf{Presentation Format:}
    \begin{itemize}
        \item Each group will present their approach and results on Friday.
        \item \textbf{Duration:} 30 minutes per group.
        \begin{itemize}
            \item $\sim$20 minutes for the presentation.
            \item 5 minutes for Q\&A.
            \item 5 minutes extra time to switch between presenting teams.
        \end{itemize}
    \end{itemize}
    
    \vspace{1em}
    \textbf{Schedule Blocks:}
    \begin{itemize}
        \item \textbf{First Block:} 10:00 -- 12:00
        \item \textbf{Second Block:} Not needed
    \end{itemize}
\end{frame}

\begin{frame}{Summary}
    \begin{itemize}
        \item \textbf{The Task:} Deblur images using Deep Learning.
        \item \textbf{The Baseline:} Total Variation regularisation solved via Condat-V\~{u} primal-dual splitting.
        \item \textbf{The Goal:} Beat the baseline using PyTorch and \texttt{deepinv}, while respecting the 100 MB and CPU runtime constraints.
    \end{itemize}
    
    \vspace{1.5em}
    \begin{center}
        \Large \textbf{Good luck, and happy coding!}
    \end{center}
\end{frame}

\begin{frame}[allowframebreaks]{References}
    \bibliography{references}
\end{frame}

\end{document}